% =============================================================================
% Capítulo 7 — Discussão
% =============================================================================
\chapter{Discussão}
\label{cap:discussao}

Este capítulo interpreta os resultados apresentados no
\autoref{cap:resultados}, analisando o comportamento exibido pelo agente PPO
em comparação com o \textit{baseline} FSM, discutindo as limitações
observadas e avaliando em que medida a proposta central do trabalho ---
utilizar proveniência como fonte de \textit{Reward Shaping} para aprendizado
por reforço --- foi validada pelos dados coletados.

% -----------------------------------------------------------------------------
\section{Análise de Comportamento e Agressividade}
\label{sec:analise-comportamento}

Os dados quantitativos da avaliação revelam um perfil comportamental
marcadamente distinto entre as duas condições experimentais. A análise a
seguir examina as dimensões mais relevantes dessa distinção.

\subsection{Duração e Persistência}
\label{subsec:duracao-persistencia}

A duração média das sessões contra o chefe PPO ($15{,}18 \pm 7{,}74$ s) é
aproximadamente 2,5 vezes superior à duração média contra o chefe FSM
($5{,}97 \pm 2{,}79$ s). Esse resultado indica que o chefe PPO é capaz de
prolongar o combate de forma consistente, resistindo por mais tempo antes de
ser derrotado. A maior variabilidade observada na condição PPO (desvio-padrão
de $7{,}74$ s, contra $2{,}79$ s no FSM) reflete a natureza menos previsível
do agente treinado: enquanto o FSM segue padrões fixos que conduzem a durações
relativamente uniformes, o PPO produz sessões com dinâmicas variadas, oscilando
entre 8,8 e 36,1 segundos na amostra oficial.

A persistência do chefe PPO não decorre de uma estratégia defensiva
elaborada --- os dados mostram que o agente é predominantemente ofensivo ---,
mas sim da pressão constante que exerce sobre o jogador. Ao atacar e se
movimentar com alta frequência, o agente dificulta a aproximação do jogador
e impõe um ritmo de combate mais intenso, o que prolonga naturalmente a
duração das sessões.

\subsection{Agressividade Ofensiva}
\label{subsec:agressividade}

A diferença mais expressiva entre as duas condições reside na contagem de
ataques do chefe. O chefe FSM executou, em média, $5{,}60 \pm 1{,}58$ ataques
por sessão, enquanto o chefe PPO executou $296{,}60 \pm 186{,}86$ --- uma
diferença de aproximadamente \textbf{53 vezes}. A distribuição de ações do
agente PPO confirma essa postura ofensiva: 39,05\% de todas as decisões do
agente correspondem à ação \textit{Attack}, seguida por \textit{Jump} com
27,95\%. Juntas, essas duas ações representam aproximadamente 67\% do
comportamento total do agente.

Essa agressividade reflete-se no comportamento adaptativo do jogador. Nas
sessões contra o PPO, o jogador executou mais pulos ($3{,}50$ vs.\ $0{,}30$),
mais ataques ($15{,}80$ vs.\ $10{,}70$) e mais \textit{dashes} ($0{,}50$ vs.\
$0{,}20$), sugerindo que a pressão exercida pelo chefe PPO forçou uma postura
mais ativa e evasiva por parte do jogador. A necessidade de mais ataques para
vencer indica que o jogador precisou investir mais esforço ofensivo para
superar o chefe PPO, ao passo que contra o FSM bastavam estratégias mais
diretas e econômicas.

\subsection{Capacidade de Ameaça}
\label{subsec:capacidade-ameaca}

O dano médio recebido pelo jogador foi ligeiramente superior na condição PPO
($56{,}00 \pm 22{,}21$) em comparação com a condição FSM ($50{,}00 \pm
14{,}14$). Embora a diferença nas médias seja moderada --- em parte limitada
pelo \textit{cooldown} de dano que equaliza a taxa de dano efetivo entre as
condições ---, o dado qualitativamente mais significativo é a ocorrência da
\textbf{única derrota do jogador} em toda a amostra de 20 sessões.

Nessa sessão específica, o jogador recebeu 100 pontos de dano (totalidade de
sua vida) e o chefe sobreviveu com 50 pontos de vida remanescentes. Essa
derrota demonstra que o agente PPO é capaz de produzir cenários de combate nos
quais o chefe representa uma ameaça real à sobrevivência do jogador, ao
contrário do FSM, cuja previsibilidade permitiu ao jogador vencer todas as
sessões sem exceção.

% -----------------------------------------------------------------------------
\section{O Fenômeno da Repetição Mecânica}
\label{sec:repeticao-mecanica}

A despeito dos indicadores positivos de agressividade e persistência, o
comportamento do agente PPO apresenta uma limitação que deve ser discutida com
transparência: a alta repetição de ações, comumente referida como
\textit{spam}.

\subsection{Caracterização do Fenômeno}
\label{subsec:caracterizacao-spam}

Os dados de distribuição de ações revelam que o agente concentra 39,05\%
de suas decisões na ação \textit{Attack} e 27,95\% na ação \textit{Jump}.
Essa concentração produz um comportamento visualmente repetitivo: o chefe PPO
ataca e pula de forma quase contínua, com pouca variação tática aparente. A
consequência visual é que o chefe pode parecer ``preso'' na animação de
ataque, pois o \textit{feedback} visual representa a última ação executada e
o agente ataca com altíssima frequência.

Embora o \textit{cooldown} de dano (0,8 segundos) implementado no
\texttt{BossAttackDamage} limite a taxa de dano efetivo --- impedindo que o
\textit{spam} resulte em morte instantânea do jogador ---, a contagem de ações
registrada no CSV reflete todas as \textit{tentativas} de ataque, não apenas
aquelas que efetivamente aplicaram dano. Assim, a grande maioria dos 296,60
ataques médios por sessão não gera dano real; eles representam decisões do
agente que são neutralizadas pelo mecanismo de \textit{cooldown}.

\subsection{Origem do Comportamento}
\label{subsec:origem-comportamento}

A repetição mecânica pode ser atribuída à combinação de dois fatores
estruturais:

\begin{enumerate}
    \item \textbf{Ausência de custo no espaço de ações}: o espaço de ações
          discretas do \texttt{BossAgent} não incorpora \textit{cooldowns}
          internos nem penalidades por repetição. Cada uma das seis ações pode
          ser executada em qualquer passo de decisão, sem restrição temporal
          ou custo energético. O algoritmo PPO, ao identificar que ações de
          ataque geram recompensa --- tanto extrínseca (por causar dano)
          quanto auxiliar (via \textit{Reward Shaping}) ---, convergiu para
          uma política que maximiza a frequência dessas ações;

    \item \textbf{Modelo congelado sem pós-processamento}: o modelo avaliado
          é o resultado direto do \textit{run} de treinamento, exportado como
          ONNX e integrado sem ajustes. A decisão de preservar o modelo
          congelado é metodologicamente correta --- a avaliação deve medir o
          comportamento aprendido, não um comportamento corrigido
          \textit{post hoc} --- mas implica que qualquer tendência degenerada
          da política aprendida é transportada integralmente para a avaliação.
\end{enumerate}

Essa observação sugere que, em trabalhos futuros, a inclusão de
\textit{cooldowns} diretamente no espaço de ações do agente ou a adição de
penalidades por repetição na função de recompensa poderiam mitigar o fenômeno
sem comprometer a autonomia do PPO.

\subsection{Preservação Metodológica}
\label{subsec:preservacao-metodologica}

É importante registrar que a decisão de não corrigir o comportamento
repetitivo antes da análise foi deliberada. Corrigir o \textit{spam} ---
por exemplo, adicionando um \textit{cooldown} forçado às ações do agente
antes da coleta --- descaracterizaria a avaliação, pois introduziria uma
modificação \textit{post hoc} que alteraria os dados coletados sem refletir
uma melhoria genuína no aprendizado. O comportamento repetitivo é um
resultado experimental do modelo congelado e, como tal, deve ser reportado
e discutido, não ocultado.

% -----------------------------------------------------------------------------
\section{Validação da Proposta: Reward Shaping por Proveniência}
\label{sec:validacao-proposta}

A questão central do trabalho é se sequências extraídas de grafos de
proveniência podem ser utilizadas como sinal auxiliar de recompensa para
orientar o treinamento de um agente PPO, produzindo um comportamento
mensuravelmente distinto do \textit{baseline} FSM --- sem recorrer a
imitação direta ou demonstrações gravadas.

\subsection{Evidências de Distinção Comportamental}
\label{subsec:evidencias-distincao}

Os dados da avaliação fornecem evidências consistentes de que o objetivo foi
alcançado. O agente PPO difere do \textit{baseline} FSM em múltiplas
dimensões mensuráveis:

\begin{itemize}
    \item \textbf{Duração}: sessões 2,5 vezes mais longas;
    \item \textbf{Volume de ações}: 53 vezes mais ataques por sessão;
    \item \textbf{Ameaça}: única condição a produzir derrota do jogador;
    \item \textbf{Distribuição de ações}: política concentrada em ações
          ofensivas, contrastando com o padrão cíclico regular do FSM;
    \item \textbf{Impacto no jogador}: aumento de pulos, ataques e
          \textit{dashes} do jogador na condição PPO, indicando adaptação a
          um oponente mais exigente.
\end{itemize}

Essas diferenças não são marginais nem atribuíveis a ruído amostral: a
magnitude das diferenças --- particularmente na contagem de ataques e na
duração --- excede largamente os desvios-padrão de ambas as condições.

\subsection{O Papel da Proveniência}
\label{subsec:papel-proveniencia}

O mecanismo de \textit{Reward Shaping} por proveniência operou conforme
projetado. As três ações que participam do \textit{Reward Shaping} ---
\textit{Attack}, \textit{Jump} e \textit{Dash} --- totalizam 77,48\% de
todas as ações executadas pelo agente, indicando que o sinal auxiliar de
recompensa exerceu influência mensurável sobre a distribuição da política
aprendida. O agente desenvolveu preferência pelas ações que eram
recompensadas quando executadas em sequências compatíveis com padrões
vitoriosos extraídos dos grafos de proveniência.

Esse resultado é significativo porque demonstra que a proveniência cumpriu
sua função de \textbf{fonte de informação causal} para o treinamento: as
sequências extraídas dos grafos --- que codificam padrões de ações associados
a desfechos favoráveis --- foram efetivamente incorporadas ao comportamento
do agente por meio do \textit{Reward Shaping}, sem a necessidade de
demonstrações diretas, arquivos \texttt{.demo} ou discriminadores adversários
(GAIL).

\subsection{Distinção em Relação à Imitação Direta}
\label{subsec:distincao-imitacao}

É fundamental distinguir o resultado obtido de um resultado de imitação. O
agente PPO \textbf{não copiou} o comportamento do jogador humano. Ele
explorou o ambiente livremente via PPO e recebeu recompensa adicional quando
executou sequências de ações compatíveis com padrões vitoriosos. O
comportamento resultante --- altamente agressivo e repetitivo --- não se
assemelha ao comportamento do jogador humano nas sessões originais de coleta
de proveniência. Isso confirma que o agente desenvolveu sua própria política
a partir dos sinais de recompensa, não uma réplica do comportamento observado.

Essa distinção é relevante tanto do ponto de vista técnico --- o pipeline não
utiliza GAIL \cite{ho2016gail} nem demonstrações gravadas --- quanto do ponto
de vista da contribuição acadêmica: o trabalho demonstra uma via alternativa
para orientar o aprendizado de agentes de jogo, baseada em informação causal
rastreável.

\subsection{Qualificações Necessárias}
\label{subsec:qualificacoes}

Os resultados devem ser interpretados com as seguintes qualificações, que
delimitam o alcance da contribuição:

\begin{enumerate}
    \item A amostra de 20 sessões (10 por condição) permite identificar
          diferenças comportamentais observadas, mas não estabelece
          significância estatística formal;
    \item A avaliação foi conduzida pelo próprio desenvolvedor, que conhece
          os padrões de ambos os chefes, introduzindo potencial viés de
          familiaridade;
    \item O modelo PPO não foi retreinado após a adição do
          \texttt{PlayerHealth} real; as observações de vida do jogador e
          \textit{cooldown} de \textit{dash} permaneceram como
          \textit{placeholders} fixos;
    \item A cena PPO não exportou grafos de proveniência durante a avaliação,
          por decisão de escopo;
    \item Os dados documentam um comportamento \textit{distinto}, não
          necessariamente \textit{superior}: as métricas provam diferença
          comportamental observada, não superioridade universal do agente PPO
          sobre o FSM.
\end{enumerate}

\subsection{Síntese}
\label{subsec:sintese-discussao}

Em síntese, o trabalho demonstrou que sequências extraídas de grafos de
proveniência podem ser utilizadas como sinal auxiliar de recompensa em um
agente PPO, produzindo um chefe com comportamento mensuravelmente distinto
do \textit{baseline} FSM em um ambiente de jogo 2D. O agente resultante é
mais agressivo, mais persistente e capaz de ameaçar a sobrevivência do
jogador, embora o comportamento aprendido apresente repetição mecânica que
compromete a percepção de sofisticação tática. A contribuição reside na
demonstração de viabilidade da integração entre proveniência e aprendizado
por reforço em um protótipo jogável, abrindo caminhos para refinamentos
futuros.

% -----------------------------------------------------------------------------
\section{Trabalhos Futuros}
\label{sec:trabalhos-futuros}

As limitações identificadas indicam direções concretas para trabalhos futuros:

\begin{enumerate}
    \item \textbf{Cooldowns no espaço de ações}: incorporar
          \textit{cooldowns} diretamente às ações discretas do agente,
          impedindo a repetição imediata e incentivando maior variedade
          tática;
    \item \textbf{Retreino com observações completas}: retreinar o modelo
          com a vida real do jogador e o \textit{cooldown} de \textit{dash}
          como observações efetivas, permitindo que o agente adapte seu
          comportamento ao estado do oponente;
    \item \textbf{\textit{Curriculum learning}}: introduzir fases de
          treinamento com dificuldade crescente, permitindo que o agente
          desenvolva estratégias progressivamente mais sofisticadas;
    \item \textbf{Avaliação com playtesters externos}: conduzir avaliação
          com jogadores externos e instrumentos de percepção subjetiva
          (questionários), complementando as métricas quantitativas com
          dados de experiência do jogador;
    \item \textbf{\textit{Self-play}}: investigar o uso de \textit{self-play}
          como mecanismo complementar de treinamento, permitindo que o agente
          evolua contra versões anteriores de si mesmo;
    \item \textbf{\textit{Reward Shaping} adaptativo}: explorar esquemas de
          \textit{Reward Shaping} que atualizem as sequências-alvo ao longo
          do treinamento, à medida que novas sessões de proveniência são
          coletadas.
\end{enumerate}
